\documentclass[11pt,letterpaper,oneside]{report}
\usepackage[T1]{fontenc}
\usepackage[utf8]{inputenc}
\usepackage{lmodern}
\usepackage[margin=0.90in,headheight=15pt]{geometry}
\usepackage{microtype}
\usepackage[svgnames,table]{xcolor}
\usepackage{amsmath,amssymb,booktabs,longtable,tabularx,array}
\usepackage{enumitem}
\usepackage[most]{tcolorbox}
\usepackage{listings}
\usepackage{fancyhdr}
\usepackage{titlesec}
\usepackage{hyperref}
\usepackage{xurl}
\usepackage{lastpage}
\definecolor{ink}{HTML}{14273C}
\definecolor{primary}{HTML}{0B5379}
\definecolor{accent}{HTML}{B66A30}
\definecolor{pale}{HTML}{F1F5F8}
\definecolor{codebg}{HTML}{F5F7F8}
\hypersetup{colorlinks=true,linkcolor=primary,urlcolor=primary,pdftitle={Linux and UNIX System Programming Engineering Suite}}
\titleformat{\chapter}[display]{\normalfont\LARGE\bfseries\color{ink}}{\color{accent}\large\MakeUppercase{\chaptername}\ \thechapter}{10pt}{}[\vspace{0.4ex}\titlerule]
\titleformat{\section}{\Large\bfseries\color{primary}}{\thesection}{0.7em}{}
\titleformat{\subsection}{\large\bfseries\color{ink}}{\thesubsection}{0.7em}{}
\setlist[itemize]{leftmargin=*,topsep=3pt,itemsep=2pt}
\setlist[enumerate]{leftmargin=*,topsep=3pt,itemsep=3pt}
\setlength{\parskip}{5pt}
\setlength{\parindent}{0pt}
\lstdefinestyle{sysc}{language=C,basicstyle=\ttfamily\footnotesize,backgroundcolor=\color{codebg},keywordstyle=\bfseries\color{primary},commentstyle=\color{ForestGreen},stringstyle=\color{Brown},breaklines=true,showstringspaces=false,tabsize=4,frame=single,rulecolor=\color{LightSteelBlue},columns=fullflexible,keepspaces=true}
\lstset{style=sysc}
\newtcolorbox{engineeringnote}[1][]{colback=pale,colframe=primary,boxrule=0.55pt,arc=1.5mm,title=Engineering note,fonttitle=\bfseries,#1}
\newtcolorbox{modernnote}[1][]{colback=white,colframe=accent,boxrule=0.55pt,arc=1.5mm,title=Modern extension --- not in the 2010 textbook,fonttitle=\bfseries,#1}
\newcommand{\sys}[1]{\texttt{#1}}
\pagestyle{fancy}\fancyhf{}\fancyhead[L]{\small\color{ink} Linux / UNIX Engineering Suite}\fancyhead[R]{\small\color{primary}\nouppercase{\leftmark}}\fancyfoot[C]{\small\thepage\ / \pageref{LastPage}}
\renewcommand{\headrulewidth}{0.45pt}
\setcounter{tocdepth}{1}
\setlength{\emergencystretch}{2em}
\begin{document}
\begin{titlepage}\centering\vspace*{2cm}{\color{primary}\Large ENGINEERING REFERENCE SERIES\par}\vspace{1.25cm}{\color{ink}\Huge\bfseries Secure Systems Programming Guide\par}\vspace{0.6cm}{\Large Threat modeling, privilege boundaries and defense in depth\par}\vspace{2cm}\rule{0.75\textwidth}{1.2pt}\par\vspace{0.55cm}{\large An original companion to Michael Kerrisk\textquotesingle s 2010 text}\par\vspace{0.3cm}{\large Edition 1.0\quad |\quad October 2026}\par\vfill{\small Document 05\quad |\quad pdfLaTeX-compatible standalone source}\par\end{titlepage}
\tableofcontents\clearpage
\chapter{Scope and threat model}
Threats cross process, file, privilege, IPC and network boundaries. This guide builds upon textbook Chapters 8--9, 15--18, 20--22, 27, 37--39, 43 and 56--63. Security advice here is engineering guidance; the original textbook predates the Linux-only mechanisms grouped in the modern extension sections.

\section{Assets, principals, entry points and decisions}
\begin{longtable}{p{.20\linewidth}p{.32\linewidth}p{.37\linewidth}}\toprule
Boundary & Attacker-controlled surface & Defensive invariant\\\midrule\endhead
Filesystem & pathnames, symlinks, mount view, permissions & Authorize an opened object, not a stale pathname\\
Launch & arguments, environment, CWD, inherited descriptors & Explicit executable and environment, CLOEXEC by default\\
Privilege & UID/GID and capability transitions & Reduce privileges before processing untrusted input\\
Memory & lengths, allocations, offsets, mapped regions & Overflow checks, bounds, lifetime and initialization\\
IPC & shared mappings, queues and local socket peers & Verify identity and protocol state, restrict object permissions\\
Network & byte streams, framing, connection counts, timeouts & Bound parsing, resource use and wait time\\
Signals & async callbacks and race-prone control flow & Signal-safe handler operations and managed waits\\
\bottomrule\end{longtable}\section{Misuse-case taxonomy}
\begin{itemize}
\item Unauthorized file read through symlink replacement or directory race
\item Command execution through shell interpolation or polluted PATH/LD\_* environment
\item Privilege retention when application code no longer needs elevated rights
\item Out-of-bounds access from untrusted lengths, integer overflow, and unsafe C string functions
\item Resource exhaustion through unbounded file descriptors, processes, threads, queues or network frames
\item Privilege confusion when process identity, container identity and peer credentials are assumed equivalent
\end{itemize}
\chapter{Files and race-resistant operations}
\section{Path resolution protocol}
\begin{enumerate}
\item Acquire a trusted directory handle where possible; document how it was established.
\item Open relative to the trusted handle, choosing O\_CLOEXEC and leaf O\_NOFOLLOW when applicable.
\item Use fstat on the descriptor to inspect the actual object; do not substitute access() plus open().
\item For stronger Linux resolution constraints evaluate openat2 with RESOLVE\_BENEATH, RESOLVE\_IN\_ROOT, or relevant restrictive flags; verify unsupported-kernel handling.
\item Apply size, type, ownership and permission rules after opening, considering hard-link and mount semantics.
\item Keep the descriptor stable for use; define cleanup and file replacement behavior.
\end{enumerate}
O\_NOFOLLOW restricts only a particular final symlink case. It does not provide a comprehensive solution for untrusted directory ancestors, magic links, mount traversals or filesystem changes. The needed protection depends on the attacker and path model.

\begin{lstlisting}
#define _POSIX_C_SOURCE 200809L
#include <errno.h>
#include <fcntl.h>
#include <stdio.h>
#include <stdlib.h>
#include <sys/stat.h>
#include <unistd.h>

/* Caller must already hold a trustworthy directory fd.
   This prevents symlinks at the final component, but not all
   possible traversal across attacker-controlled parents. */
static int read_regular_leaf(int dirfd, const char *leaf) {
    int fd = openat(dirfd, leaf, O_RDONLY | O_CLOEXEC | O_NOFOLLOW);
    if (fd < 0) return -1;
    struct stat st;
    if (fstat(fd, &st) < 0) {
        int saved = errno;
        close(fd);
        errno = saved;
        return -1;
    }
    if (!S_ISREG(st.st_mode)) {
        close(fd);
        errno = EINVAL;
        return -1;
    }
    char c;
    ssize_t n = read(fd, &c, 1);
    int saved = errno;
    close(fd);
    errno = saved;
    return n >= 0 ? 0 : -1;
}
int main(void) {
    int dir = open(".", O_RDONLY | O_DIRECTORY | O_CLOEXEC);
    if (dir < 0) { perror("open dir"); return 1; }
    int rc = read_regular_leaf(dir, "input.txt");
    if (rc < 0) perror("read_regular_leaf");
    close(dir); return rc < 0;
}
\end{lstlisting}
\section{Temporary and durable files}
\begin{itemize}
\item Avoid constructing predictable names and avoid world-writable scratch locations without a deliberate trust model.
\item Use mkstemp or suitably constrained platform APIs, restrictive initial modes, exclusive creation and directory-fd operations.
\item Treat fsync of file data and directory entries separately when crash-consistency is an explicit requirement.
\item Do not blindly delete a pathname after concurrent attacker-controlled replacement; prefer scoped ownership of temporary directories.
\end{itemize}
\chapter{Privileged process lifecycle}
Use a narrowly privileged broker, early privilege reduction, or capabilities when the platform and threat model permit. Distinguish real, effective and saved IDs and supplementary groups. Dropping effective UID alone can leave a recoverable privilege path; the exact call and order matter.

\section{Launching untrusted work}
\begin{itemize}
\item Do not use system(), popen(), or shell-string construction on untrusted data.
\item Specify absolute executable path where practical, sanitized environment and explicit argv.
\item Set close-on-exec at acquisition time to avoid races with concurrent fork and exec.
\item Apply process limits and file-descriptor closures before work starts.
\item Log process exit statuses, timeouts and resource failures without exposing secrets.
\end{itemize}
\section{Privilege review checklist}
\begin{enumerate}
\item Why is elevation needed at all? Is a broker or system service interface available?
\item What exactly can the child open, create, connect to and execute?
\item Can the process regain IDs or capabilities after dropping them?
\item Are dynamic libraries and plugin search paths trusted?
\item How are failures reported without revealing sensitive credentials or content?
\end{enumerate}
\chapter{Signals, concurrency and memory safety}
Async signal handlers may invoke only functions permitted by the applicable standard. A robust application typically blocks designated signals in all threads and receives them in a dedicated thread using sigwait; use self-pipe or signal descriptors where appropriate. Ordinary shared variables are not synchronization just because a signal handler modifies them.

\begin{itemize}
\item Use pthread mutexes and condition variables to protect invariants; do not use unsynchronized volatile state as a substitute for atomic or locking guarantees.
\item Validate length multiplication before allocating: count > SIZE\_MAX / element\_size indicates overflow.
\item Clean secrets from memory where possible, but recognize compiler optimization and swap or core-dump limitations; use platform primitives where applicable.
\item Treat mmap regions as shared, private or file-backed according to explicit mapping semantics; avoid sharing process pointers across remapped address spaces.
\item Constrain queue capacities and thread counts to prevent denial of service.
\end{itemize}
\chapter{Modern Linux containment layers}
\section{no\_new\_privs}
Prevent certain future privilege gains across exec; often a prerequisite for unprivileged seccomp filter installation.

\begin{modernnote}
No substitute for per-resource authorization.
\end{modernnote}
\section{seccomp-bpf}
Filter system call entry points using kernel mediation; test architecture, permitted actions, and exact syscall numbers.

\begin{modernnote}
Syscall filtering does not understand higher-level file or network policy; wrapper calls can change.
\end{modernnote}
\section{Landlock}
Confine selected file-system and other accessible operations via a ruleset subject to ABI support.

\begin{modernnote}
ABI-version checks and inherited rules require careful setup.
\end{modernnote}
\section{Namespaces}
Change views of PID, user, mount, network and IPC resources, often used in container implementations.

\begin{modernnote}
Not automatically a security barrier; user-namespace rules and mount composition are delicate.
\end{modernnote}
\section{cgroup v2}
Enforce hierarchical CPU, memory, I/O and PID limits to reduce resource-exhaustion blast radius.

\begin{modernnote}
A resource controller does not enforce per-operation trust decisions.
\end{modernnote}
\section{Linux capabilities}
Grant subsets of historically root-gated operations; reason about bounding, inheritable, ambient, permitted and effective sets.

\begin{modernnote}
Capabilities are privileges, not sandbox features by themselves.
\end{modernnote}
\chapter{Network attack surface and observability}
\begin{itemize}
\item Define accepted address families and bind addresses; demos bind loopback to avoid exposing unauthenticated services.
\item Implement maximum frame sizes and budgets for concurrent connections, unread bytes and pending writes.
\item Use a vetted TLS implementation for encrypted and authenticated network transport; raw sockets do not provide TLS.
\item Restrict sensitive logs: include request IDs, timing, errno categories and peer context without secrets.
\item Test malformed packets, truncated frames, slow clients, EOF during headers and repeated connection churn.
\end{itemize}
\chapter{Security assurance worksheet}
\begin{longtable}{p{.18\linewidth}p{.73\linewidth}}\toprule
Control & Objective verification\\\midrule\endhead
Path safety & Symlink swap race harness, untrusted parent-directory test, alternate mount scenario\\
Privileges & Before/after credential and capability snapshots, negative permission checks\\
Process launch & Fuzzed argv/env and executable failure tests, file descriptor inheritance inspection\\
Resource bounds & File count, thread count, memory and CPU stress with graceful failure\\
Memory & ASan/UBSan, hardening flags and input corpus regression\\
Protocol & Truncation, invalid length, replay/out-of-order and timeout cases\\
Containment & Allowed behavior still works; blocked behavior fails as intended on target kernel\\
\bottomrule\end{longtable}\section{Example risk register}
\begin{longtable}{p{.12\linewidth}p{.29\linewidth}p{.26\linewidth}p{.22\linewidth}}\toprule
ID & Failure mode & Mitigation & Proof\\\midrule\endhead
SYS-01 & Path-swap read & Dirfd anchoring and controlled resolution & Race harness\\
SYS-02 & Child execution injection & argv/env, no shell & Adversarial arguments\\
SYS-03 & Unbounded connections & Per-peer and global quotas & Load test\\
SYS-04 & Privilege survives startup & Permanent drop + verify & Credential assertions\\
SYS-05 & Data corruption on crash & Sync/rename protocol & Crash tests on target FS\\
\bottomrule\end{longtable}
\chapter*{Source and currency notes}
\addcontentsline{toc}{chapter}{Source and currency notes}
\textbf{Foundational source:} Michael Kerrisk, \emph{The Linux Programming Interface: A Linux and UNIX System Programming Handbook}, No Starch Press, 2010, ISBN 978-1-59327-220-3. The textbook's chapter titles and start pages organize the study guide; explanations, tasks, and code in this suite are independently authored educational material, not reproductions of the original text or its example programs.

\textbf{Current reference baseline:} The Open Group / IEEE, \emph{POSIX.1-2024, Issue 8}, \url{https://pubs.opengroup.org/onlinepubs/9799919799/}; Linux man-pages project, \url{https://man7.org/linux/man-pages/}; Linux kernel documentation, \url{https://docs.kernel.org/}. For Linux-only extensions consult the relevant per-interface manual pages and kernel documentation for the target kernel, architecture and libc; availability can differ materially.

\textbf{Scope:} This is an engineering study and implementation companion. It does not claim to restate a normative standard or guarantee a specific Linux kernel's behavior. System calls, libc wrappers and higher-level library functions are explicitly distinguished where material.
\end{document}
